\section{Conclusion}\label{sec:conclusion}

\begin{proof}[Completion of Proposition~\ref{prop:anchor}]
The residual cases treated above are exhaustive.  If \(A[2]\) is
reducible, its invariant line contains a nonzero rational two-torsion
point, and Section~\ref{sec:reducible} supplies the anchor.  Otherwise
the image is an irreducible subgroup of
\(\operatorname{GL}_2(\mathbb F_2)\simeq S_3\), hence is \(C_3\) or
\(S_3\).  The \(S_3\) case divides according to whether its quadratic
sign field is imaginary or real; these are handled in
Sections~\ref{alternation-imaginary:anchor-section} and~\ref{sec:real}.
Proposition~\ref{cyclic:anchor} handles \(C_3\).

Each construction gives an odd fundamental discriminant \(h\),
allowing one, prime to \(2N_A\), and an imaginary fundamental
discriminant \(k\) prime to \(2hN_A\), with every prime dividing
\(2hN_A\) split in the corresponding quadratic field.  It first
proves a simple zero of the product
\(L(A^h,s)L(A^{hk},s)\), and then proves
\(X(A^h)+X(A^{hk})=0\) with both discrepancies defined.  These are
all the assertions of Proposition~\ref{prop:anchor}.
\end{proof}

The Positive comparison, Proposition~\ref{prop:positive}, was
completed in Section~\ref{sec:positive-completion}; the CM comparison,
Proposition~\ref{prop:cm}, was completed in Section~\ref{sec:cm}.
We can therefore finish the proof for the original curve.

\begin{proof}[Proof of Theorem~\ref{thm:main}]
Let \(E/\mathbb Q\) satisfy \(s_2(E)\le1\).  Suppose first that
\(E\) is non-CM.  Proposition~\ref{prop:anchor} gives \(h,k\) with
\[
 \an(E^h)+\an(E^{hk})=1,
 \qquad X(E^h)+X(E^{hk})=0.
\]
Both analytic orders are nonnegative integers, so each is at most
one.  By the first part of Proposition~\ref{prop:positive}, both
terms in the second equality are nonnegative.  Hence
\[
 X(E^h)=X(E^{hk})=0.
\]
The discriminants \(h,k\) are odd and coprime.  Their product is
again a fundamental discriminant, and exactly one of \(h,hk\) is
positive because \(k<0\).  Denote that positive member by \(a\),
allowing \(a=1\).  It is prime to \(2N_E\), and
\(\an(E^a)\le1\), \(X(E^a)=0\).  The second part of
Proposition~\ref{prop:positive}, applied to the original curve \(E\),
therefore gives
\[
 \an(E)=s_2(E),\qquad X(E)=0.
\]
For a CM curve, Proposition~\ref{prop:cm} gives the same conclusion
directly.

In either case the analytic order is zero or one.  The classical
Gross--Zagier--Kolyvagin theorem gives
\[
 \operatorname{rank}E(\mathbb Q)=\an(E)=s_2(E),
 \qquad \#\Sha(E/\mathbb Q)<\infty
\]
for the whole Tate--Shafarevich group, not just its two-primary
part~\cite{GZ86,Kolyvagin90}.  The period rationality in rank zero,
and the split Gross--Zagier and arithmetic-volume comparisons in rank
one, give \(Q_E\in\mathbb Q_{>0}\), as established in
Section~\ref{sec:classical}.  By the definition of the discrepancy,
\[
 0=X(E)=v_2\!\left(\frac{Q_E}{\#\Sha(E/\mathbb Q)}\right),
\]
which is precisely
\(v_2(Q_E)=v_2(\#\Sha(E/\mathbb Q))\).

Throughout, \(\Omega_E\) is the period over the whole real group,
and \(\Reg_E\) is computed on the full Mordell--Weil lattice with
the height pairing specified in Section~\ref{sec:introduction}.  Thus the equality
contains the real-component contribution and every point-index
square required by the stated normalization.  The residual
classification and the dyadic local comparisons cover every residual
representation and every reduction type at two.
\end{proof}


\begin{corollary}[Exact two-primary BSD for almost all quadratic twists]
\label{cor:quadratic-twist-density}
Fix an elliptic curve \(E/\mathbb Q\). For each nonzero squarefree integer
\(d\), write \(E^d\) for the quadratic twist of \(E\) by \(d\), and put
\[
 \begin{aligned}
 \mathcal D(X)&=\{d\in\mathbb Z:0<|d|\le X,\ d\text{ squarefree}\},\\
 \mathcal D_j(E;X)&=\{d\in\mathcal D(X):s_2(E^d)=j\}
 \qquad(j=0,1).
 \end{aligned}
\]
Then
\[
 \lim_{X\to\infty}\frac{\#\mathcal D_j(E;X)}{\#\mathcal D(X)}
 =\frac12\qquad(j=0,1).
\]
For every nonzero squarefree \(d\) with \(s_2(E^d)=j\in\{0,1\}\), put
\(F=E^d\). Then
\[
 \operatorname{rank}F(\mathbb Q)
 =\operatorname{ord}_{s=1}L(F,s)=s_2(F)=j,
 \qquad \#\Sha(F/\mathbb Q)<\infty,
\]
and
\[
 Q_F=\frac{L^{(j)}(F,1)\,(\#T_F)^2}
 {j!\,\Omega_F\Reg_F\prod_{\ell\ \mathrm{finite}}c_\ell(F)}
 \in\mathbb Q_{>0},
 \qquad v_2(Q_F)=v_2\bigl(\#\Sha(F/\mathbb Q)\bigr).
\]
Here the period, regulator, torsion, and Tamagawa factors have the
normalizations fixed before Theorem~\ref{thm:main}. In particular, the
exact two-primary formula holds on the density-one union of these
subfamilies, whose common ranks are respectively zero and one. The
density is in the full signed squarefree family for fixed \(E\); the
formula asserts equality only at the prime two.
\end{corollary}

\begin{proof}
The proof of Theorem~1.2 in \cite{OpenAIGoldfeldDensity2026} gives density
\(1/2\) for each of the usual full \(2\)-power Selmer coranks zero and
one among the signed squarefree parameters ordered by absolute value.
Its \(c_2\) is \(s_2\) here, since both use the usual local Kummer
conditions at every place. This gives the two displayed density limits.
For every twist in either subfamily, Theorem~\ref{thm:main} applied to
\(F=E^d\) gives all the pointwise assertions. The two disjoint
subfamilies have densities summing to one.
\end{proof}
